% =====================================================================
%  Chương V — KẾT LUẬN VÀ HƯỚNG PHÁT TRIỂN
%  File này được thiết kế để \input vào file luận văn chính.
%  Không yêu cầu gói đặc biệt ngoài các gói đã dùng ở các chương trước.
% =====================================================================
\chapter{V. KẾT LUẬN VÀ HƯỚNG PHÁT TRIỂN}
\label{chap:ketluan}

% =====================================================================
\section{5.1. Kết luận}
\label{sec:kl_ketluan}

Khóa luận đã xây dựng và đánh giá một hệ thống trích xuất bộ ba
\textit{(aspect term, aspect category, sentiment)} cho đánh giá khách sạn
tiếng Việt, tổ chức thành bốn stage (tiền xử lý $\to$ ATE bằng T5 $\to$
APC bằng \texttt{FastLcfBertMultiTask} $\to$ tổng hợp bộ ba), với hai
phương pháp kỹ thuật trọng tâm là \textbf{Local Context Focus (LCF)} và
\textbf{Token Merging (ToMe)}. Các kết quả chính:

\begin{enumerate}
  \item \textbf{Kiến trúc Joint vượt Pipeline khoảng 18 điểm Micro F1}
        ($\approx 71$ so với $\approx 53$). Chỉ số Oracle Joint Acc gần
        như không đổi giữa hai kiến trúc cho thấy nút thắt nằm ở bước
        trích xuất khía cạnh (ATE), không phải bộ phân loại -- một minh
        chứng rõ ràng cho hiện tượng lan truyền lỗi.
  \item \textbf{Mô hình tốt nhất}: Joint $\cdot$ T5 $\cdot$
        \texttt{seq\_resize} đạt Micro F1 $= 71.50$ và Oracle Joint
        $= 89.10$; Joint $\cdot$ BERT $\cdot$ \texttt{lcf\_seq\_cdm\_resize}
        đạt Micro F1 $= 71.18$ với chi phí huấn luyện thấp hơn.
  \item \textbf{CDM tốt hơn CDW trên nhãn thiểu số}: mặt nạ nhị phân CDM
        nâng Macro F1 thêm $6$--$9$ điểm so với trọng số mềm CDW, trong
        khi Micro F1 gần như không đổi.
  \item \textbf{Gộp token không làm giảm chất lượng}: các cấu hình ToMe
        đạt độ chính xác ngang hoặc nhỉnh hơn baseline, khẳng định việc
        loại token dư thừa không phá hỏng biểu diễn cần cho phân loại.
\end{enumerate}

\subsection{5.1.1. Về quyết định chỉ dùng chế độ resize}
\label{subsec:kl_resize}

Toàn bộ thực nghiệm sử dụng chế độ \texttt{resize} (nội suy chuỗi đã gộp
trở lại độ dài gốc). Đây là một lựa chọn \emph{có chủ đích} nhằm cô lập
\textbf{câu hỏi chất lượng biểu diễn}: giữ độ dài và kiến trúc các lớp sau
không đổi, khác biệt duy nhất giữa các cấu hình là tác động của thao tác
gộp. Khóa luận do đó trả lời được câu hỏi ``gộp token theo chiến lược nào
thì bảo toàn độ chính xác'' một cách sạch sẽ.

Hệ quả cần nêu rõ một cách trung thực: vì ToMe được đặt \emph{sau}
backbone và chuỗi được nội suy lại về $L = 128$, \textbf{chế độ resize
không mang lại tăng tốc suy luận} -- thực tế còn làm tăng nhẹ thời gian do
chi phí gộp và nội suy. Việc hiện thực hóa lợi ích tốc độ của Token
Merging là nội dung của hướng phát triển dưới đây.

% =====================================================================
\section{5.2. Hạn chế}
\label{sec:kl_hanche}

\begin{itemize}
  \item \textbf{Chưa đo tăng tốc thực sự.} Với cấu hình resize hiện tại,
        đóng góp của ToMe nằm ở khía cạnh chất lượng biểu diễn, chưa phải
        ở hiệu năng tính toán như mục tiêu ban đầu hướng tới.
  \item \textbf{Nút thắt ATE.} ATE F1 Pipeline chỉ $\approx 60$ do lỗi
        biên cụm từ và chọn nhầm khía cạnh, giới hạn hiệu suất đầu-cuối.
  \item \textbf{Mất cân bằng dữ liệu.} Một số tổ hợp nhãn
        (\texttt{BRANDING\_neutral}, \texttt{FACILITY\_neutral}) đạt
        F1 $= 0$ do quá ít mẫu; weighted loss và dữ liệu bổ trợ chỉ giảm
        nhẹ vấn đề.
  \item \textbf{Quy mô đánh giá.} Bộ dữ liệu tương đối nhỏ (test 312 mẫu)
        và thuộc một miền (khách sạn), nên khả năng khái quát ra miền khác
        chưa được kiểm chứng.
\end{itemize}

% =====================================================================
\section{5.3. Hướng phát triển}
\label{sec:kl_huongphattrien}

\subsection{5.3.1. Hiện thực hóa tăng tốc của Token Merging}

Đây là hướng quan trọng nhất, nối tiếp trực tiếp phân tích ở
mục~\ref{subsec:kl_resize}:

\begin{itemize}
  \item \textbf{Chế độ compact (\texttt{resize=False}).} Giữ nguyên chuỗi
        rút gọn $L'$ để các lớp self-attention phía sau xử lý ít token
        hơn, rồi đo trực tiếp độ trễ/thông lượng so với baseline.
  \item \textbf{Pre-BERT ToMe (\texttt{use\_pre\_tome}).} Gộp token ngay ở
        mức embedding \emph{trước} 12 lớp Transformer -- đòn bẩy tốc độ
        lớn nhất, vì rút ngắn chính phần tốn kém $O(n^2)$ của backbone.
        Mã nguồn đã hỗ trợ sẵn (\texttt{\_pre\_bert\_merge}) nhưng chưa
        đưa vào bộ cấu hình chính.
  \item \textbf{Báo cáo đánh đổi tốc độ--độ chính xác} (latency, FLOPs,
        throughput so với Micro/Macro F1) để định lượng lợi ích thực sự
        của từng chiến lược gộp.
\end{itemize}

\subsection{5.3.2. Cải thiện bước ATE}

Vì ATE là nút thắt, có thể: tinh chỉnh T5 với ràng buộc biên chặt hơn,
dùng decoding có ràng buộc (constrained decoding) để đầu ra luôn là chuỗi
con hợp lệ, hoặc thử kiến trúc trích xuất theo span thay cho sinh tự do.

\subsection{5.3.3. Xử lý mất cân bằng và mở rộng dữ liệu}

Bổ sung dữ liệu cho các nhãn hiếm, thử các kỹ thuật như focal loss,
over-sampling có kiểm soát, hoặc sinh dữ liệu tăng cường; đồng thời mở
rộng đánh giá sang miền khác (nhà hàng, thương mại điện tử) để kiểm chứng
khả năng khái quát.

\subsection{5.3.4. Khảo sát sâu hơn về LCF và chiến lược gộp}

Quét siêu tham số ngưỡng SRD $\alpha$ và số bước gộp; kết hợp tín hiệu
attention thực (thay vì trọng số đồng nhất) cho chiến lược
\texttt{attention\_weighted}; và thử các chiến lược gộp mới nhận biết
ranh giới khía cạnh tốt hơn.

% =====================================================================
\section*{Lời kết}
\addcontentsline{toc}{section}{Lời kết}

Khóa luận đã cung cấp một so sánh hệ thống giữa các chiến lược gộp token
khi kết hợp với cơ chế LCF trong một mô hình ABSA đa nhiệm, trên dữ liệu
đánh giá khách sạn tiếng Việt. Kết quả cho thấy gộp token bảo toàn được
chất lượng phân loại và CDM có lợi cho nhãn thiểu số, đồng thời chỉ ra rõ
ràng con đường để biến tiềm năng giảm chi phí của Token Merging thành lợi
ích tăng tốc đo được trong các nghiên cứu tiếp theo.
